\chapter{Graph Minors}

\section{Well-quasi-ordering}

\begin{lemma}
  \label{QO_tricolor} \lean{QO_tricolor} \uses{theo:ramsey_infgraph} \leanok
  If \(X\) is a quasi-ordered set, in any infinite sequence, there exists either an infinite increasing subsequence, an infinite strictly decreasing subsequence or an infinite antichain
\end{lemma}

\begin{proof}
  Let \(\leq\) be a quasi-ordering on \(X\)

  Let \((x_i)_{i\in\mathbb{N}}\) be an infinite sequence in \(X\)

  Let \(K\) be the complete graph on \(\mathbb{N}\)

  Let \(C=\{0,1,2\}\) a set of colors

  For each edge \((i,j)\in E(K)\), if \(i<j\) we assign to it the color \(0\) if \(x_i\leq x_j\), \(1\) if \(x_i>x_j\) and \(2\) if \(x_i\) and \(x_j\) are incomparable

  By Ramsey's Theorem, there exists an subset \(S\subset E(K)\) which is infinite and monochromatic

  Therefore, if its color is \(0\), then for all \((i,j)\in S, x_i\leq x_j\), if its color is \(1\), then for all \((i,j)\in S, x_i>x_j\) and if its color is \(2\), then for all \((i,j)\in S\), \(x_i\) and \(x_j\) are incomparable

  It follows that there exists either an infinite increasing subsequence, an infinite strictly decreasing subsequence or an infinite antichain
\end{proof}

\begin{proposition}
  \label{WQO_iff} \lean{WQO_iff} \leanok \mathlibok
  A quasi-ordering $\leq$ on $X$ is a well-quasi-ordering if and only if $X$
  contains neither an infinite antichain nor an infinite strictly decreasing
  sequence $x_0 > x_1 > \cdots$.
\end{proposition}

\begin{proof}
  \leanok
  Diestel's book uses Ramsey theory to prove this but in fact it is essentially
  in mathlib already.
\end{proof}

\begin{corollary}
  \label{WQO_incsubseq} \uses{WQO_iff}
  If $X$ is well-quasi-ordered, then every infinite sequence in $X$ has an
  infinite increasing subsequence.
\end{corollary}

% \begin{proof}
%   Let \(\leq\) be a quasi-order on \(X\). 
  
%   Suppose that \(X\) is well-quasi-ordered.

%   Let \((x_i)_{i\in\mathbb{N}}\) an infinite sequence in \(X\).

%   There exists an index \(i\) such that \(\{j>i/x_i\leq x_j\}\) is infinite.

%   Indeed, if we suppose not, for all index \(i\), there exists \(j>i\) such that \(x_i\not\leq x_j\)

%   Therefore, as \((x_i)_{i\in\mathbb{N}}\) is infinite, we can construct an infinite subsequence \((x_{i_k})_{k\in\mathbb{N}}\) such that \(\forall k\in\mathbb{N}, x_{i_k}>\x_{i_{k+1}}\) or \(x_{i_k}\) and \(\x_{i_{k+1}}\) are incomparable
  
%   Suppose there is an infinite subsequence of incomparable elements, this is an antichain which contradicts the well-quasi-ordered property.

%   If not, there exists an index \(k_*\) such that \(\forall k>k_*, x_{i_k}>\x_{i_{k+1}}\), this is an infinite strictly decreasing sequence which also contradicts the well-quasi-ordered property.

%   Let \(i_0\) be an index such that \(\{j>i_0/x_{i_0}\leq x_j\}\) is infinite.

%   We proceed by induction. 
  
%   Let \(n\in\mathbb{N}\), we suppose constructed \(x_{i_0},...,x_{i_n}\) such that \(i_0<...<i_n\), \(x_{i_0}\leq ...\leq x_{i_n}\) and \(A_n:=\{j>i_n/x_{i_n}\leq x_j\}\) is infinite.
  
%   \((x_i)_{i\in A_n}\) is an infinite sequence in \(X\), therefore, according to the previous result, there exists an index \(i_{n+1}\in A_n\) such that \(A_{n+1}=\{j>i_{n+1}/x_{i_{n+1}}\leq x_j\}\) is infinite.
  
%   As \(i_{n+1}\in A_n\), it follows that \(i_n<i_{n+1}\) and \(x_{i_n}\leq x_{i_{n+1}}\)
  
%   We have obtained \(x_{i_0},...,x_{i_{n+1}}\) such that \(i_0<...<i_{n+1}\), \(x_{i_0}\leq ...\leq x_{i_{n+1}}\) and \(A_{n+1}=\{j>i_n/x_{i_{n+1}}\leq x_j\}\) is infinite, which concludes the induction
  
%   We have constructed an infinite increasing subsequence \((x_{i_k})_{k\in\mathbb{N}}\) of \((x_i)_{i\in\mathbb{N}}\)
% \end{proof}

\begin{theorem}[Higman's Lemma, Diestel 12.1.3]
  \label{theo:Higman} \lean{Higman} \leanok \uses{WQO_incsubseq}
  If $X$ is well-quasi-ordered by, then so is $[X]^{<ω}$.
\end{theorem}

\begin{proof}
  Let \(\leq\) be a well-quasi-ordering on \(X\),

  Suppose that \(\leq\) is not a well-quasi-ordering on \([X]^{<\omega}\)

  For all \(A,B\in[X]^{<\omega}\), we write \(A\leq B\) if there exists \(f:A\to B\) injective such that for all \(a\in A\), \(a\leq f(a)\)

  A sequence \(\left(B_i\right)_{i\in\mathbb{N}}\in([X]^{<\omega})^\mathbb{N}\) is bad if for all \(i,j\in\mathbb{N}\) with \(i<j\), \(B_i\not\leq B_j\)

  A finite sequence \(A_0,...,A_n\) is a prefix of \(\left(B_i\right)_{i\in\mathbb{N}}\in([X]^{<\omega})^\mathbb{N}\) if for all \(i\leq n\), \(A_i=B_i\)

  \(A_0,...,A_n\in[X]^{<\omega}\) is a bad prefix if there exists a bad sequence \(\left(A_i\right)_{i\in\mathbb{N}}\in([X]^{<\omega})^\mathbb{N}\) such that \(T_0,...,T_n\) is a prefix of \(\left(A_i\right)_{i\in\mathbb{N}}\in([X]^{<\omega})^\mathbb{N}\)

  \(A\in[X]^{<\omega}\) is minimal in \(\mathcal{A}\subset[X]^{<\omega}\) if for all \(A'\in\mathcal{A}, |A|\leq |A'|\)

  We proceed by induction with the following hypothesis. For any given \(n\in\mathbb{N}\), for all bad prefix \(A_0,...,A_{n-1}\in[X]^{<\omega}\), there exists \(A_n\) minimal in \(\{A'/A_0,...,A_{n-1},A_n\text{ is a bad prefix}\}\)

  As \(\leq\) is not a well-quasi-ordering on \([X]^{<\omega}\), there exists \(\left(B_i\right)_{i\in\mathbb{N}}\in([X]^{<\omega})^\mathbb{N}\) a bad sequence.

  Let \(\mathcal{B}=\{B\in[X]^{<\omega}/B\text{ is a bad prefix}\}\) and \(\mathcal{B}_\#=\{n\in\mathbb{N}/\exists B\in\mathcal{B},|B|=n\}\)

  \(B_0\) is a bad prefix so \(\mathcal{B}\) is non-empty so \(\mathcal{B}_\#\) is also non-empty

  \(\mathcal{B}_\#\) is a non-empty subset of \(\mathbb{N}\) so it has a minimum \(m\in\mathbb{N}\)

  There exists \(A_0\in\mathcal{B}\) such that \(|A_0|=m\)

  \(A_0\) is minimal in \(\mathcal{B}\) which proves the base case

  Let \(n\in\mathbb{N}\), 

  Let \(A_0,...,A_{n-1}\in[X]^{<\omega}\) be a bad prefix

  Let \(\mathcal{B}=\{B\in[X]^{<\omega}/A_0,...,A_{n-1},B\text{ is a bad prefix}\}\) and \(\mathcal{B}_\#=\{n\in\mathbb{N}/\exists B\in\mathcal{B},|B|=n\}\)

  \(\mathcal{B}\) is non-empty so \(\mathcal{B}_\#\) is too

  \(\mathcal{B}_\#\) is a non-empty subset
  
  of \(\mathbb{N}\) so it has a minimum \(m\in\mathbb{N}\)

  There exists \(A_n\in\mathcal{B}\) such that \(|A_n|=m\)

  \(A_n\) is minimal in \(\mathcal{B}\)

  The sequence \(\left(A_n\right)_{n\in\mathbb{N}}\) is bad in \([X]^{<\omega}\) by construction
  
  In particular, \(A_n\neq\emptyset\) as for all \(A\in[X]^{<\omega}, \emptyset\leq A\)

  For all \(n\in\mathbb{N}\), we choose \(a_n\in A_n\)

  Let \(C_n:=A_n\backslash\{a_n\}\)

  \(X\) is well-quasi-ordered, therefore, according to the corollary 12.1.2, the sequence has an infinite increasing subsequence \(\left(A_{n_i}\right)_{i\in\mathbb{N}}\)

  The sequence \(A_0,A_1,...,A_{n_0-1},C_{n_0},C_{n_1},C_{n_2},...\) is good

  That fact is true because \(|C_{n_0}|<|A_{n_0}|\) and \(|A_{n_0}|\) is minimal

  There exists a good pair in the sequence. It can have the form \((A_i,A_j)\), \((A_i,B_j)\) or \((B_i,B_j)\) with \(i<j\)

  The first case is impossible as \(\left(A_n\right)_{n\in\mathbb{N}}\) is bad

  The second case is also impossible as \(B_j\leq A_j\)

  Thus, it has the form \((B_i,B_j)\)

  There exists \(f:B_i\to B_j\) injective such that for all \(a\in B_i\), \(a\leq f(a)\)

  \(a_i\leq a_j\) as \(\left(A_{n_i}\right)_{i\in\mathbb{N}}\) is an infinite increasing sequence

  Therefore, we can extend the injection \(f\) by setting \(f(a_i)=a_j\)

  We obtain that \((A_i,A_j)\) is good which contradicts the hypothesis
\end{proof}

\section{The graph minor theorem for trees}

\begin{theorem}[Kruskal 1960]
  \label{theo:Kruskal} \lean{Kruskal} \uses{Highman}
  The finite trees are well-quasi-ordered by the topological minor relation \(\preceq\)
\end{theorem}

\begin{proof}
  \(\left(A_i\right)_{i\in\mathbb{N}}\) is bad if \(\forall i,j\in\mathbb{N}, i<j\Rightarrow A_i\npreceq A_j\) and is good if it is not bad

  \(T_0,...,T_n\) is a prefix of \(\left(A_i\right)_{i\in\mathbb{N}}\) if \(\forall i\leq n, A_i=T_i\)

  \(T_0,...,T_n\) is a bad prefix if it is a prefix of some bad sequence

  \(T\) is minimal in \(\mathcal{T}\) if \(\forall T'\in\mathcal{T}, T'\preceq T\Rightarrow T\preceq T'\)

  For any tree \(T\) with root \(r\), for all \(u,v\in T\), we write \(u<_Tv\) if the unique path from \(r\) to \(v\) passes through \(u\). \(<_T\) is the tree-order of \(T\)

  If \(T=(V,E,r)\) is a rooted tree, we define the forest \(T-r=(V\backslash\{r\},E',R)\) with \(E'=\{(v,v')\in E/v,v'\neq r\}\) and \(R\) the set of all children of \(r\) in \(T\)

  Let \(T\) be a tree and \(x,y\in T\), if \(x\) is an ascendant of \(y\), we denote the unique path from \(x\) to \(y\) by \(x[\to]y\)
  
  By contradiction, suppose that finite trees are not well-quasi-ordered by \(\preceq\)

  There exists \(\left(A_i\right)_{i\in\mathbb{N}}\) a bad sequence

  By induction, we show that for any given \(n\in\mathbb{N}\), for all bad prefix \(T_0,...,T_{n-1}\) the set \(\mathcal{T}=\{T/T_0,...,T_{n-1},T\text{ is a bad prefix}\}\) is non-empty and there exists \(T_n\) minimal in \(\mathcal{T}\)
  
  Let \(T_0,...,T_{n-1}\) be a bad prefix, and \(\mathcal{T}=\{T/T_0,...,T_{n-1},T\text{ is a bad prefix}\}\)

  \(T_0,...,T_{n-1}\) is a bad prefix so there exists \(\left(A'_i\right)_{i\in\mathbb{N}}\) a bad sequence with \(T_0,...,T_{n-1}\) as a prefix

  \(A'_n\in\mathcal{T}\) so \(\mathcal{T}\) is non-empty

  For all \(T\in\mathcal{T}\), \(\empty\preceq T\) and \(\preceq\) is a preorder

  Therefore, by Zorn's lemma for preorders, there exists \(T_n\) minimal in \(\mathcal{T}\)

  The base case is shown the same way with an empty bad prefix and with \(A_0\in\mathcal{T}\)

  The sequence \(\left(T_i\right)_{i\in\mathbb{N}}\) is bad by construction

  For each \(n\in\mathbb{N}\), we denote the root of \(T_n\) by \(r_n\)

  For each \(n\in\mathbb{N}\), let \(B_n\) be the forest \(T_n-r_n\)

  Let \(B=\displaystyle\bigcup_{n\in\mathbb{N}}B_n\)

  Let \(\left(T^{(k)}\right)_{k\in\mathbb{N}}\)

  For all \(k\in\mathbb{N}\), there exists \(n_k\in\mathbb{N}\) such that \(T^{(k)}\in\B_{n_k}\)

  Let \(N=\min\{n_k\in\mathbb{N}/k\in\mathbb{N}\}\) and \(m=\min\{k\in\mathbb{N}/N=n_k\}\)

  Let \(C=\left(C_i\right)_{i\in\mathbb{N}}\) with \(C_i=T_i\) if \(i<N\) and \(C_i=T^{(i-N+m)}\) if \(i\geq N\)

  \(T^{(m)}\subsetneq T_N\) and \(T_N\) is minimal in \(\mathcal{T}=\{T/T_0,...,T_{N-1},T\text{ is a bad prefix}\}\)

  Therefore \(T_0,...,T_{N-1}\) is not a bad prefix and \(C\) is good

  There exists \(i,j\in\mathbb{N}\) such that \(i<j\) and \(C_i\preceq C_j\)

  If \(i<N\) and \(j<N\), we get \(T_i\preceq T_j\) which contradicts the fact that \(\left(T_i\right)_{i\in\mathbb{N}}\) is bad

  If \(i<N\) and \(j\geq N\), we get \(T_i\preceq T^{(j-N+m)}\) but \(T^{(j-N+m)}\preceq T_{n_{j-N+m}}\) so, by transitivity, \(T_i\preceq T_{n_{j-N+m}}\)

  As \(i<N\leq j-N+m\) by minimality of \(N\) it is impossible

  If \(i\geq N\) and \(j<N\), we get \(j<i\) which is impossible as \(i<j\)

  So the good pair is in \(\left(T^{(k)}\right)_{k\in\mathbb{N}}\)

  \(B\) is well-quasi-ordered

  According to Higman's Lemma, \([B]^{<\omega}\) is also well-quasi-ordered

  Therefore, \(\left(B_n\right)_{n\in\mathbb{N}}\) has a good pair

  There exists \(i,j\in\mathbb{N}\) and \(f:B_i\to B_j\) such that \(i<j\), \(f\) injective and for all \(T\in\B_i, T\preceq f(T)\)

  For all \(T\in B_i\) there exists an embedding \(\varphi_T:T\to f(T)\) that preserves the tree-order of \(T\)

  Let \(\varphi:T_i\to T_j\) be a function such that \(\forall T\in B_i, \varphi_{|T}=\varphi_T\) and \(\varphi(r_i)=r_j\) for every \(v\in T_i\) child of \(r_i\), \(\varphi(r_iv)=r_j[\to]\varphi(v)\)

  \(\varphi\) preserves the tree-order of \(T_i\), therefore \(T_i\preceq T_j\) which contradicts the fact that \(\left(T_i\right)_{i\in\mathbb{N}}\) is bad
\end{proof}

\setcounter{section}{6}
\section{The graph minor theorem}

\begin{theorem}
  \label{theo:GraphMinorTheorem} \lean{GraphMinorTheorem} \uses{def:minor} \leanok
  The finite graphs are well-quasi-ordered by the minor relation $\preceq$.
\end{theorem}

